\section{Strojové učení. Druhy úloh strojového učení, klasifikace algoritmů, dostupné nástroje}


\subsection*{Vymezení strojového učení}

Strojové učení je oblast umělé inteligence, ve které se systém učí ze vzorových dat a na jejich základě
zlepšuje své rozhodování, predikce nebo rozpoznávání vzorů. Rozdíl proti klasickému programování je v tom,
že pravidla nejsou plně ručně naprogramována, ale jsou odhadnuta z dat.

Typicky máme:
\[
X = (x_1,\ldots,x_p)
\]
jako vstupní atributy, prediktory nebo features, a u učení s učitelem také cílovou proměnnou:
\[
y.
\]

Cílem modelu je naučit se funkci:
\[
\hat{f}: X \rightarrow y,
\]
která dobře funguje nejen na trénovacích datech, ale hlavně na nových datech.

\paragraph{Intuice.}
U klasického programování řekneme počítači přesná pravidla. U strojového učení mu dáme příklady a algoritmus
se pokusí pravidla najít sám. Například místo ruční definice, co je spam, model dostane mnoho označených
e-mailů a učí se, jaké vlastnosti se spamem souvisí.

\subsection*{Pattern recognition a data mining}

V materiálech se objevuje pojem \term{pattern recognition}, tedy rozpoznávání vzorů. Na vstupu jsou surová
data a na výstupu je akce, klasifikace nebo doporučení podle nalezených vzorů.

Příklady:
\begin{itemize}
    \item detekce podvodné platby,
    \item rozpoznávání řeči nebo obrazu,
    \item klasifikace zákazníků,
    \item predikce spotřeby elektřiny,
    \item analýza clickstreamů v e-shopu,
    \item text mining.
\end{itemize}

\term{Data mining} neboli dobývání znalostí z databází je proces hledání netriviálních, potenciálně
užitečných a srozumitelných vztahů v rozsáhlých datech. V praxi je silně propojen s datovými sklady,
datovou kvalitou, business znalostí a interpretací výsledků.

\subsection*{Základní druhy úloh strojového učení}

\subsubsection*{Učení s učitelem}

Učení s učitelem má k dispozici trénovací data:
\[
\{(x_i,y_i)\}_{i=1}^n.
\]

U každého objektu známe cílovou proměnnou $y_i$. Model se učí vztah mezi vstupy $x_i$ a výstupem $y_i$.

\paragraph{Klasifikace.}
Cílová proměnná je kategoriální:
\[
y_i \in \{1,\ldots,K\}.
\]

Příklady:
\begin{itemize}
    \item schválit / neschválit úvěr,
    \item spam / ne-spam,
    \item zákazník odejde / neodejde,
    \item diagnóza pacienta,
    \item segment zákazníka.
\end{itemize}

\paragraph{Regrese.}
Cílová proměnná je číselná:
\[
y_i \in \mathbb{R}.
\]

Příklady:
\begin{itemize}
    \item predikce ceny nemovitosti,
    \item predikce spotřeby energie,
    \item odhad tržeb,
    \item predikce pojistných nákladů.
\end{itemize}

\subsubsection*{Učení bez učitele}

Učení bez učitele nemá cílovou proměnnou. Máme pouze:
\[
\{x_i\}_{i=1}^n.
\]

Cílem je najít strukturu v datech.

Typické úlohy:
\begin{itemize}
    \item \term{shlukování} -- hledání skupin podobných objektů,
    \item \term{asociační pravidla} -- hledání souvislostí mezi položkami,
    \item \term{redukce dimenze} -- PCA, autoenkodéry,
    \item \term{detekce anomálií} -- hledání neobvyklých pozorování.
\end{itemize}

\paragraph{Příklad.}
U zákazníků neznáme předem segmenty. Shlukovací algoritmus se pokusí najít skupiny zákazníků podle podobného
chování, nákupů nebo demografie.

\subsubsection*{Doporučování}

Doporučovací systémy predikují, jaké položky by uživatele mohly zajímat. V materiálech se objevuje
\term{collaborative filtering}, například doporučování filmů podle matice uživatel--položka.

Typická vstupní matice:
\[
R =
\begin{pmatrix}
r_{11} & r_{12} & \cdots & r_{1m}\\
r_{21} & r_{22} & \cdots & r_{2m}\\
\vdots & \vdots & \ddots & \vdots\\
r_{n1} & r_{n2} & \cdots & r_{nm}
\end{pmatrix},
\]
kde $r_{ij}$ je hodnocení položky $j$ uživatelem $i$.

Jedna možnost je rozklad:
\[
R \approx U \Sigma V^\top,
\]
tedy singulární rozklad, který hledá latentní faktory uživatelů a položek.

\subsubsection*{Detekce anomálií}

Detekce anomálií hledá pozorování, která se výrazně liší od běžných dat. Může být:
\begin{itemize}
    \item supervised, pokud máme označené podvody/anomálie,
    \item unsupervised, pokud označení nemáme.
\end{itemize}

Příklad z praxe je detekce podvodných plateb. Pokud není k dispozici dostatek označených podvodů, lze použít
například \term{Isolation Forest}. Ten izoluje neobvyklé body pomocí náhodných stromů; anomálie bývají
izolovány menším počtem dělení.

\subsection*{Klasifikace algoritmů}

\subsubsection*{Podle typu úlohy}

\begin{itemize}
    \item \textbf{Klasifikační algoritmy:}
    rozhodovací stromy, logistická regrese, k-NN, Naive Bayes, SVM, random forest, boosting,
    neuronové sítě.
    
    \item \textbf{Regresní algoritmy:}
    lineární regrese, regresní stromy, random forest regressor, gradient boosting regressor,
    support vector regression, neuronové sítě.
    
    \item \textbf{Shlukovací algoritmy:}
    k-means, hierarchické shlukování, DBSCAN, aglomerativní metody.
    
    \item \textbf{Asociační pravidla:}
    Apriori, FP-Growth, GUHA, 4FT-Miner.
    
    \item \textbf{Redukce dimenze:}
    PCA, SVD, t-SNE, UMAP, autoenkodéry.
\end{itemize}

\subsubsection*{Podle interpretovatelnosti}

\begin{itemize}
    \item \textbf{White-box modely:}
    lineární/logistická regrese, malé rozhodovací stromy, pravidla, rozhodovací tabulky.
    
    \item \textbf{Black-box modely:}
    random forest, gradient boosting, SVM s nelineárním kernelem, neuronové sítě.
\end{itemize}

White-box model je srozumitelnější, ale nemusí mít nejlepší predikční výkon. Black-box model může být přesnější,
ale často vyžaduje dodatečné vysvětlovací metody, například LIME, SHAP nebo parciální závislosti.

\paragraph{No-Free-Lunch princip.}
Neexistuje jeden algoritmus, který by byl nejlepší pro všechny problémy. Algoritmus se vybírá podle cíle,
povahy dat, velikosti datasetu, požadavků na interpretaci, nákladů chyb a výpočetních omezení.

\subsection*{Vybrané algoritmy}

\subsubsection*{Rozhodovací strom}

Rozhodovací strom rozděluje data pomocí pravidel typu:
\[
x_j \leq c.
\]

Výsledkem je strom, jehož listy obsahují predikci třídy nebo hodnoty. Výhodou je vysoká interpretovatelnost.
Nevýhodou je sklon k přeučení.

\subsubsection*{k-nearest neighbours}

k-NN klasifikuje nový objekt podle nejbližších trénovacích objektů. Pro nový objekt $x$ najdeme $k$ nejbližších
sousedů a zvolíme nejčastější třídu:
\[
\hat{y}(x)=\operatorname{mode}\{y_i: x_i \in N_k(x)\}.
\]

Typická vzdálenost pro numerická data:
\[
d(x,z)=\sqrt{\sum_{j=1}^p (x_j-z_j)^2}.
\]

Výhody:
\begin{itemize}
    \item jednoduchý princip,
    \item žádný silný předpoklad o tvaru rozhodovací hranice,
    \item přirozená lokální interpretace přes podobné příklady.
\end{itemize}

Nevýhody:
\begin{itemize}
    \item citlivost na škálování proměnných,
    \item vysoká paměťová a výpočetní náročnost při predikci,
    \item curse of dimensionality: ve vysoké dimenzi se vzdálenosti stávají méně informativní.
\end{itemize}

\subsubsection*{Naive Bayes}

Naive Bayes vychází z Bayesovy věty:
\[
P(c|x)=\frac{P(x|c)P(c)}{P(x)}.
\]

Pro klasifikaci stačí maximalizovat čitatel:
\[
\hat{c}
=
\underset{c}{\operatorname{argmax}}
\,
P(c)P(x|c).
\]

Protože $x=(x_1,\ldots,x_p)$ může mít mnoho atributů, Naive Bayes zavádí zjednodušující předpoklad
podmíněné nezávislosti:
\[
P(x|c)=\prod_{j=1}^p P(x_j|c).
\]

Potom:
\[
\hat{c}
=
\underset{c}{\operatorname{argmax}}
\,
P(c)\prod_{j=1}^p P(x_j|c).
\]

Je „naivní“, protože atributy obvykle nezávislé nejsou. Přesto často funguje dobře, hlavně v textové
klasifikaci.

\paragraph{Laplace correction.}
Pokud se některá kombinace ve trénovacích datech nevyskytne, pravděpodobnost by byla nulová. Proto se používá
vyhlazení:
\[
\hat{P}(x_j=a|c)=
\frac{n_{a,c}+\alpha}{n_c+\alpha m},
\]
kde $m$ je počet možných hodnot atributu a $\alpha$ je obvykle 1.

\subsubsection*{SVM}

Support Vector Machine hledá nadrovinu, která odděluje třídy a maximalizuje margin:
\[
w^\top x + b = 0.
\]

U lineárně separovatelných dat:
\[
y_i(w^\top x_i+b)\geq 1.
\]

Optimalizační úloha hard-margin SVM:
\[
\min_{w,b}\frac{1}{2}\|w\|^2
\quad
\text{s.t.}
\quad
y_i(w^\top x_i+b)\geq 1.
\]

Šířka marginu je:
\[
\frac{2}{\|w\|}.
\]

U neseparovatelných dat se používá soft-margin SVM se slack proměnnými:
\[
\min_{w,b,\xi}
\frac{1}{2}\|w\|^2 + C\sum_{i=1}^n \xi_i
\]
\[
y_i(w^\top x_i+b)\geq 1-\xi_i,
\qquad
\xi_i\geq 0.
\]

Parametr $C$ řídí kompromis mezi širokým marginem a chybami klasifikace.

Pro nelineární data lze použít kernel:
\[
K(x_i,x_j).
\]

Příklady:
\[
K(x,z)=x^\top z
\quad \text{lineární kernel},
\]
\[
K(x,z)=(x^\top z+r)^d
\quad \text{polynomiální kernel},
\]
\[
K(x,z)=\exp(-\gamma\|x-z\|^2)
\quad \text{RBF kernel}.
\]

\subsubsection*{Neuronové sítě}

Neuronové sítě skládají mnoho jednoduchých výpočtových jednotek. U mělké sítě máme vstupní vrstvu, jednu
skrytou vrstvu a výstupní vrstvu. Detailněji jsou popsány v otázce 15.

\subsection*{CRISP-DM jako metodika analytického projektu}

CRISP-DM je standardní metodika pro data mining projekty. Má šest fází:

\begin{enumerate}
    \item \term{Business understanding}:
    pochopení problému, cíle, omezení, nákladů chyb a očekávaného využití.
    
    \item \term{Data understanding}:
    získání dat, pochopení atributů, kvality, rozsahu a omezení.
    
    \item \term{Data preparation}:
    čištění, spojování dat, imputace, odstranění duplicit, transformace, feature engineering.
    
    \item \term{Modeling}:
    volba algoritmů, trénování modelů, ladění hyperparametrů.
    
    \item \term{Evaluation}:
    ověření kvality modelu, kontrola, zda model řeší původní business problém.
    
    \item \term{Deployment}:
    nasazení modelu, reporting, API, automatizace, monitoring a údržba.
\end{enumerate}

V praxi bývá nejvíce práce v komunikaci s vlastníkem dat, čištění, přípravě dat, interpretaci a zajištění
důvěry ve výsledky, ne v samotném spuštění algoritmu.

\subsection*{Příprava dat}

Před modelováním je nutné řešit:

\begin{itemize}
    \item \textbf{Chybějící hodnoty:}
    odstranění řádků, imputace průměrem/mediánem, regresní imputace, modelová imputace.
    
    \item \textbf{Škálování:}
    důležité pro k-NN, SVM, neuronové sítě a clustering.
    
    \item \textbf{Kódování kategorií:}
    dummy proměnné, one-hot encoding, ordinal encoding.
    
    \item \textbf{Feature engineering:}
    tvorba nových atributů z doménové znalosti.
    
    \item \textbf{Feature selection:}
    výběr relevantních proměnných.
    
    \item \textbf{Data leakage:}
    nesmí se použít informace, která by v okamžiku predikce nebyla dostupná.
\end{itemize}

Standardizace:
\[
z_{ij}=\frac{x_{ij}-\bar{x}_j}{s_j}.
\]

Min-max škálování:
\[
x_{ij}^{*}=
\frac{x_{ij}-\min(x_j)}
{\max(x_j)-\min(x_j)}.
\]

\subsection*{Rozdělení dat a validace}

Dostupná označená data se obvykle rozdělí na:

\begin{itemize}
    \item \term{trénovací množinu} -- model se na ní učí,
    \item \term{validační množinu} -- ladění hyperparametrů a výběr modelu,
    \item \term{testovací množinu} -- finální nezávislé ověření kvality.
\end{itemize}

Jednoduché rozdělení:
\[
D = D_{\text{train}} \cup D_{\text{test}}.
\]

\subsubsection*{Křížová validace}

U $K$-fold cross-validation se data rozdělí na $K$ částí. Model se $K$krát natrénuje, vždy na $K-1$ částech,
a ověří na zbývající části.

Průměrná chyba:
\[
CV_K
=
\frac{1}{K}
\sum_{k=1}^K Err_k.
\]

\paragraph{Stratifikace.}
U nevyvážených tříd je vhodné zachovat podíly tříd v každém foldu.

\paragraph{Time-based split.}
U časově závislých dat se nesmí náhodně míchat minulost a budoucnost. Trénujeme na minulosti a testujeme
na pozdějších datech.

\subsection*{Hodnocení klasifikačních modelů}

Pro binární klasifikaci:

\[
\begin{array}{c|cc}
 & \text{Skutečně pozitivní} & \text{Skutečně negativní}\\
\hline
\text{Predikováno pozitivní} & TP & FP\\
\text{Predikováno negativní} & FN & TN
\end{array}
\]

\subsubsection*{Accuracy}

\[
Accuracy=
\frac{TP+TN}{TP+TN+FP+FN}.
\]

Accuracy může být zavádějící u nevyvážených tříd.

\subsubsection*{Precision}

\[
Precision=
\frac{TP}{TP+FP}.
\]

Říká, jaká část pozitivních predikcí byla skutečně pozitivní.

\subsubsection*{Recall / sensitivity}

\[
Recall=
\frac{TP}{TP+FN}.
\]

Říká, jaká část skutečně pozitivních případů byla nalezena.

\subsubsection*{Specificity}

\[
Specificity=
\frac{TN}{TN+FP}.
\]

Říká, jaká část skutečně negativních případů byla správně označena jako negativní.

\subsubsection*{F1 score}

\[
F_1=
2\cdot
\frac{Precision\cdot Recall}
{Precision+Recall}.
\]

F1 je vhodné, pokud chceme vyvážit precision a recall.

\subsubsection*{ROC a AUC}

ROC křivka zobrazuje vztah:
\[
TPR = \frac{TP}{TP+FN}
\]
proti:
\[
FPR = \frac{FP}{FP+TN}.
\]

AUC shrnuje ROC křivku do jednoho čísla:
\[
AUC=1
\]
znamená perfektní klasifikátor,
\[
AUC=0.5
\]
odpovídá náhodnému klasifikátoru.

\subsection*{Cost-sensitive evaluation}

Ne všechny chyby mají stejný náklad. Například v medicíně může být falešně negativní chyba mnohem
dražší než falešně pozitivní.

Celkový náklad:
\[
Cost
=
C_{TP}TP+C_{FP}FP+C_{FN}FN+C_{TN}TN.
\]

Klasifikační práh lze zvolit tak, aby minimalizoval náklad:
\[
t^*
=
\underset{t}{\operatorname{argmin}}
\,
Cost(t).
\]

\subsection*{Overfitting a underfitting}

\term{Underfitting} znamená, že model je příliš jednoduchý a špatně zachytí strukturu dat.
Má vysokou chybu na trénovacích i testovacích datech.

\term{Overfitting} znamená, že model se naučil i šum a náhodné zvláštnosti trénovací množiny.
Má nízkou chybu na trénovacích datech, ale vysokou chybu na testovacích datech.

\[
Err_{\text{train}} \ll Err_{\text{test}}
\quad \Rightarrow \quad
\text{podezření na overfitting}.
\]

Řešení:
\begin{itemize}
    \item jednodušší model,
    \item regularizace,
    \item pruning stromů,
    \item více dat,
    \item cross-validation,
    \item ensemble metody,
    \item správné ladění hyperparametrů.
\end{itemize}

\subsection*{Ladění hyperparametrů}

Parametry se učí z dat. Hyperparametry nastavujeme před učením.

Příklady:
\begin{itemize}
    \item počet sousedů $k$ u k-NN,
    \item hloubka stromu,
    \item počet stromů v random forest,
    \item learning rate u boostingu nebo neuronových sítí,
    \item parametr $C$ a kernel u SVM.
\end{itemize}

Metody ladění:
\begin{itemize}
    \item grid search,
    \item random search,
    \item cross-validation,
    \item nested cross-validation.
\end{itemize}

\subsection*{Dostupné nástroje}

\begin{itemize}
    \item \textbf{Python:}
    pandas, numpy, scikit-learn, matplotlib, seaborn, scipy, xgboost, lightgbm, tensorflow, pytorch.
    
    \item \textbf{R:}
    caret, tidymodels, randomForest, e1071, nnet, rpart.
    
    \item \textbf{GUI nástroje:}
    Weka, RapidMiner, KNIME, IBM SPSS Modeler, SAS Enterprise Miner.
    
    \item \textbf{Akademické a specializované nástroje:}
    LISp-Miner, GUHA, 4FT-Miner.
    
    \item \textbf{Cloud a produkční nástroje:}
    BigML, cloudové ML služby, kontejnery, API, CI/CD, monitoring modelů.
\end{itemize}

Materiály zmiňují práci v Pythonu a scikit-learn, cloudové nástroje, BigML pro stromy a LISp-Miner jako
systém navazující na GUHA.

\subsection*{Shrnutí otázky}

Strojové učení je přístup, kdy se model učí ze vzorových dat. Základní dělení je učení s učitelem,
učení bez učitele, doporučování a detekce anomálií. Učení s učitelem zahrnuje klasifikaci a regresi,
učení bez učitele shlukování, asociační pravidla a redukci dimenze.

Typický ML projekt se řídí logikou CRISP-DM: pochopení businessu, pochopení dat, příprava dat, modelování,
vyhodnocení a nasazení. Kvalita modelu se nehodnotí na trénovacích datech, ale na validačních/testovacích
datech pomocí vhodných metrik. Je nutné hlídat overfitting, data leakage, nevyvážené třídy a náklady chyb.

\newpage